\scanpage{26}\mspage{25}
We start all over!

\uline{Definition-scholium} A \uline{country} \inserted{$C$} is defined by a set $\mathcal{L}$ (the set of \uline{loci} of the country), and the datum of certain sets of subsets of $\mathcal{L}$, called the \uline{figures} \inserted{(of the country)}.

The set $\mathcal{F}$ of figures of $C$ is therefore a \deleted{set} subset of $\mathcal{P}(\mathcal{P}(\mathcal{L}))$ (thus $F\in \mathcal{P}(\mathcal{P}(\mathcal{P}(\mathcal{L})))$). If \deleted{$\unsure{\mathcal{F}|F}$}$F$ is a figure, its elements\textsuperscript{$\times$/} (which are subsets of $\mathcal{L}$) are called the \uline{strata} (\uline{closed}) of the figure $F$. \unsure{There|Lá|The} union of the strata \deleted{closed} of $F$ is called the \uline{support} of $F$, \inserted{denoted $|F|$.} We lay down \uline{axioms of countries}, which will be brought out later.

\uline{Cont}~1 \quad a) The empty subset of $\mathcal{P}(\mathcal{L})$ is a \deleted{stratum} figure, called ``\uline{empty figure}''. It is therefore the unique figure that has no strata.

% slanted
\margin{Contained in \uline{Cont} 5 below}
\margin{$\big\Vert$}
b) $\forall\,x\in \mathcal{L}$, \deleted{the}\inserted{there is a} figure having \deleted{the} unique stratum $\{x\}$. We call it the \uline{point figure} defined by $x$
